%!TEX TS-program = xelatex
%!TEX encoding = UTF-8 Unicode
% =====================================================================
% resume.master.tex -- MASTER RÉSUMÉ
% ---------------------------------------------------------------------
% One source for every application. Nothing here is sent as-is:
%   scripts/build.sh master   -> build/master/resume-master-catalog.pdf
%                                every bullet, with its [id | role types]
%   scripts/build.sh bases    -> build/bases/resume-<type>.pdf
%   scripts/build.sh app applications/<file>.tex
%                             -> build/applications/<file>/Carlos-Tejada-Resume.pdf
% Compiling this file directly (e.g. in VS Code) gives the catalog.
%
% Role types (archetypes):
%   ds  = Data Scientist            mle = ML / AI Engineer
%   de  = Data Engineer             rs  = Research Scientist (HCI, accessibility)
%   uxr = Quantitative / mixed-methods UX Researcher
%
% How the content is organized
%   * Each \bul{id}{role types}{text} is ONE accomplishment in ONE phrasing.
%     Variants of the same accomplishment share a stem (uw-system, uw-system.ds).
%   * Order inside a role = print order, strongest first.
%   * Numbers, titles and dates come from master/facts.tex (\Fact... macros).
%   * To add something new: add a \bul with a new id under the right role,
%     tag it, rebuild the catalog. Never add a fact only in an application file.
% =====================================================================
\documentclass[11pt, letterpaper]{awesome-resume}
% =====================================================================
% master/preamble.tex
% Shared layout and machinery for resume.master.tex and
% cover-letter.master.tex. Loaded right after \documentclass.
% You shouldn't need to edit this file to apply for a job.
% =====================================================================

% ---------------------------------------------------------------------
% 1. Page and text layout (tested fixes from the Résumé Audit, 2026-09-11)
% ---------------------------------------------------------------------
\geometry{letterpaper, left=1.5cm, right=1.5cm, top=1.1cm, bottom=1.2cm, footskip=.5cm}
\fontdir[fonts/]
\definecolor{awesome}{HTML}{222A68}
\definecolor{flagcolor}{HTML}{D35400}
\setbool{acvSectionColorHighlight}{false}
\XeTeXinterwordspaceshaping=2                 % real spaces in the PDF text (pypdf reads words, not "Datascientistwith...")
%% ...but shaping=2 also merges consecutive \href annotations on one line: the
%% second link's start is swallowed, so the first link's hotspot stretches over
%% the whole line. Every header link pointed at LinkedIn because of this, and
%% tel:/mailto: produced no annotation at all. An \mbox around each link keeps
%% its annotation intact. Use \slink, never a bare \href, wherever two or more
%% links share a line.
\newcommand{\slink}[2]{\mbox{\href{#1}{#2}}}
\hyphenpenalty=10000 \exhyphenpenalty=10000   % never split a word across lines
\setlength{\parskip}{0pt}
\setlength{\parindent}{0pt}
\pagestyle{empty}                              % no footer: parsers read footers as body text
\raggedright
\usepackage{needspace}                        % keep a role's header with its first bullet

% ---------------------------------------------------------------------
% 2. Build modes
%    Catalog  : \archetype = all  -> every bullet, with its ID and tags
%    Base     : \archetype = ds | mle | de | rs | uxr
%    Final    : application builds (\FinalBuildFlag defined) -> no highlights,
%               but every unverified fact / unfilled slot is logged and the
%               build script refuses to finish.
% ---------------------------------------------------------------------
\providecommand{\archetype}{all}
\newif\ifcatalog
\ifdefstring{\archetype}{all}{\catalogtrue}{\catalogfalse}
\newif\ifFinalBuild
\ifdefined\FinalBuildFlag \FinalBuildtrue \else \FinalBuildfalse \fi

% ---------------------------------------------------------------------
% One-page mode: \def\OnePageFlag{1} in an application file (or
% scripts/build.sh bases --onepage). Prints the same facts, less of them:
% each role keeps only its first \onepagemax bullets, roles wrapped in
% \unlessonepage are dropped, and publications shrink to the summary line.
% A bullet an application force-includes is never dropped.
% ---------------------------------------------------------------------
\newif\ifonepage
\ifdefined\OnePageFlag \onepagetrue \else \onepagefalse \fi
\newcounter{bulincluded}
\def\OnePageMax{2}
\newcommand{\onepagemax}[1]{\def\OnePageMax{#1}}
\newcommand{\unlessonepage}[1]{\ifonepage\else#1\fi}

% \unverified[label]{value}: a fact that disagrees between old files or needs your records
\newcommand{\unverified}[2][]{%
  \ifFinalBuild #2\else\textcolor{flagcolor}{#2\textsuperscript{?}}\fi
  \typeout{UNVERIFIED FACT: #1 = \detokenize{#2}}}
% \slot{what to write}: a placeholder that must be filled before sending
\newcommand{\slot}[1]{%
  \textcolor{red}{\textbf{[[SLOT: #1]]}}%
  \typeout{UNFILLED SLOT: \detokenize{#1}}}

% ---------------------------------------------------------------------
% 3. Role-type (archetype) helpers
% ---------------------------------------------------------------------
\newcommand{\archcheck}[1]{\ifdefstring{\archetype}{#1}{\def\archhit{1}}{}}
% \ifarch{ds,mle}{if the current role type is listed}{otherwise}
\newcommand{\ifarch}[3]{\def\archhit{0}\forcsvlist{\archcheck}{#1}\ifdefstring{\archhit}{1}{#2}{#3}}
% \onlyfor{rs,uxr}{content}: printed for those role types and in the catalog
\newcommand{\onlyfor}[2]{\ifcatalog #2\else\ifarch{#1}{#2}{}\fi}
% \catalogonly{content}: printed only in the catalog
\newcommand{\catalogonly}[1]{\ifcatalog #1\fi}

% Per-application include/exclude lists (set in an application file)
\providecommand{\ExcludeIDs}{}
\providecommand{\IncludeIDs}{}
\def\excludedidslist{}
\def\includedidslist{}
\renewcommand*{\do}[1]{\listadd\excludedidslist{#1}}
\expandafter\docsvlist\expandafter{\ExcludeIDs}
\renewcommand*{\do}[1]{\listadd\includedidslist{#1}}
\expandafter\docsvlist\expandafter{\IncludeIDs}

% ---------------------------------------------------------------------
% 4. Bullets: \bul{id}{role types}{text} ... \printbullets
%    A bullet prints when the current role type is in its list (or in the
%    catalog), unless the application excludes its ID; an application can
%    force-include any ID. An application can also reword one bullet with
%    \expandafter\def\csname override@<id>\endcsname{new text}.
% ---------------------------------------------------------------------
\newcommand{\bulletbuffer}{}
\newcommand{\bulcheckfor}[1]{\ifdefstring{\archetype}{#1}{\def\bulhit{1}}{}}
\newcommand{\catalogtag}[2]{\ifcatalog\hspace{0.35em}{\fontsize{7.5pt}{1em}\bodyfont\color{gray}[#1\,\textbar\,#2]}\fi}
\newcommand{\bul}[3]{%
  \def\bulhit{0}%
  \ifcatalog\def\bulhit{1}\else\forcsvlist{\bulcheckfor}{#2}\fi
  \ifinlist{#1}{\excludedidslist}{\def\bulhit{0}}{}%
  \ifonepage
    \ifdefstring{\bulhit}{1}%
      {\ifnum\value{bulincluded}<\OnePageMax\relax\else\def\bulhit{0}\fi}{}%
  \fi
  \ifinlist{#1}{\includedidslist}{\def\bulhit{1}}{}%
  \ifdefstring{\bulhit}{1}{%
    \stepcounter{bulincluded}%
    \ifcsdef{override@#1}%
      {\appto\bulletbuffer{\item \csuse{override@#1}\catalogtag{#1}{#2}}}%
      {\appto\bulletbuffer{\item #3\catalogtag{#1}{#2}}}}{}%
}
\newenvironment{rolebullets}{%
  \begin{itemize}[leftmargin=1.05em, labelsep=0.45em, itemsep=1.2pt, parsep=0pt, topsep=1.5pt, partopsep=0pt, label={\textbullet}]%
  \raggedright\fontsize{10pt}{12.4pt}\bodyfont\color{text}}%
  {\end{itemize}}
\newcommand{\printbullets}{%
  \ifdefempty{\bulletbuffer}{}{\begin{rolebullets}\bulletbuffer\end{rolebullets}}%
  \gdef\bulletbuffer{}\setcounter{bulincluded}{0}\def\OnePageMax{2}}

% ---------------------------------------------------------------------
% 5. Layout macros: header, sections, roles, skills, publications
% ---------------------------------------------------------------------
\def\theheadline{}
\def\thesummary{}
\def\theskills{}
\newcommand{\catalogentry}[2]{\par\noindent{\fontsize{8.5pt}{1em}\bodyfont\bfseries\color{awesome}[#1]}\hspace{0.4em}#2\par\vspace{3pt}}
% Headline (line under your name), summary and skills: one per role type
\newcommand{\headlinefor}[2]{\ifarch{#1}{\gdef\theheadline{#2}}{}}
\newcommand{\summaryfor}[2]{\ifcatalog\catalogentry{#1}{#2}\else\ifarch{#1}{\gdef\thesummary{#2}}{}\fi}
\newcommand{\skillsfor}[2]{\ifcatalog\catalogentry{#1}{}#2\vspace{2pt}\else\ifarch{#1}{\gdef\theskills{#2}}{}\fi}
\newcommand{\printheadline}{%
  \ifdefined\HeadlineOverride\HeadlineOverride
  \else\ifcatalog\ifdefined\CatalogHeadline\CatalogHeadline\else MASTER CATALOG \textbar{} every bullet for every role type \textbar{} not for sending\fi%
  \else\theheadline\fi\fi}
\newcommand{\printsummary}{%
  \ifcatalog\else{\fontsize{10pt}{12.6pt}\bodyfont\color{text}%
    \ifdefined\SummaryOverride\SummaryOverride\else\thesummary\fi\par}\fi}
\newcommand{\printskills}{\ifcatalog\else\ifdefined\SkillsOverride\SkillsOverride\else\theskills\fi\fi}

\newcommand{\masterheader}{%
  \begin{center}
    {\fontsize{25pt}{25pt}\headerfont\bfseries\color{darktext}\FactName\FactPostNominal}\par\vspace{2.2mm}
    {\fontsize{10.5pt}{13pt}\bodyfont\color{awesome}\printheadline}\par\vspace{1.4mm}
    {\fontsize{9.5pt}{12pt}\bodyfont\color{text}\FactCity\ifdefined\FactPhone\ \textbar\ \FactPhone\fi\ \textbar\ \FactEmail}\par\vspace{0.5mm}
    {\fontsize{9.5pt}{12pt}\bodyfont\color{text}\FactLinkedIn\ \textbar\ \FactGitHub\ \textbar\ \FactWebsite\onlyfor{rs,uxr}{\ \textbar\ \FactScholar}}\par
  \end{center}\vspace{-1mm}}

\newcommand{\msection}[1]{%
  \par\ifonepage\vspace{2pt}\else\vspace{6pt}\fi
  \noindent{\fontsize{12.5pt}{14pt}\bodyfont\bfseries\color{awesome}#1}\par\ifonepage\vspace{1pt}\else\vspace{1.5pt}\fi
  \noindent{\color{lightgray}\rule{\linewidth}{0.6pt}}\par\ifonepage\vspace{1.5pt}\else\vspace{2pt}\fi}

% \cvrole{employer}{location}{title}{dates}  -- then \bul ... \printbullets
\newcommand{\cvrole}[4]{%
  \par\needspace{5\baselineskip}\ifonepage\vspace{2.5pt}\else\vspace{4.5pt}\fi
  \noindent{\fontsize{10.5pt}{13pt}\bodyfont\bfseries\color{darktext}#1}\hfill{\fontsize{9.5pt}{13pt}\bodyfont\color{darktext}#2}\par\nopagebreak
  \noindent{\fontsize{10pt}{12.5pt}\bodyfont\color{darktext}#3}\hfill{\fontsize{9.5pt}{12.5pt}\bodyfont\color{darktext}#4}\par\nopagebreak}

% A second (earlier) appointment with the same employer: prints title + dates
% under the employer already printed by \cvrole. Harvard Law's resume guidance:
% "If you have held more than one job with the same employer, enter both jobs
% under one header." Keeps one employer/one location for a parser to read.
\newcommand{\cvsubrole}[2]{%
  \par\needspace{3\baselineskip}%
  \noindent{\fontsize{10pt}{12.5pt}\bodyfont\color{darktext}#1}\hfill{\fontsize{9.5pt}{12.5pt}\bodyfont\color{darktext}#2}\par\nopagebreak}

\newcommand{\skillline}[2]{\par\noindent\hangindent=1.2em{\fontsize{10pt}{12.6pt}\bodyfont\color{text}\textbf{#1:} #2}\par\vspace{1.5pt}}
\newcommand{\pubitem}[1]{\par\noindent\hangindent=1.2em{\fontsize{9.5pt}{12pt}\bodyfont\color{text}#1}\par\vspace{2pt}}
\newcommand{\eduline}[3]{\par\noindent{\fontsize{10pt}{12.6pt}\bodyfont\color{text}\textbf{#1}, #2}\hfill{\fontsize{9.5pt}{12.6pt}\bodyfont\color{darktext}#3}\par\vspace{1pt}}
\newcommand{\me}{\textbf{C.~E.~Tejada}}

% ---------------------------------------------------------------------
% 6. Cover-letter paragraph bank: \letterpara{id}{role types}{text}, \usepara{id}
% ---------------------------------------------------------------------
\def\paralist{}
\newcommand{\letterpara}[3]{\csgdef{para@#1}{#3}\csgdef{paratags@#1}{#2}\listgadd\paralist{#1}}
\newcommand{\usepara}[1]{\ifcsdef{para@#1}{\csuse{para@#1}}{\slot{missing paragraph "#1"}}}

% =====================================================================
% master/facts.tex -- the single source of truth
% ---------------------------------------------------------------------
% Every date, title, count and contact detail in the résumé and the
% cover letter comes from here. Tailor an application by CHOOSING and
% ORDERING content, never by changing a value in an application file.
%
% \unverified[label]{value} marks a value that disagrees between your old files,
% or one I couldn't confirm. Draft builds show it in orange with a "?".
% Application builds refuse to finish while any remain.
% When you've checked a value against your records, delete the
% \unverified{ } wrapper around it (keep the value).
% =====================================================================

% ---------------- Identity and contact ----------------
\newcommand{\FactName}{Carlos E. Tejada}
% Post-nominal on the resume and letter header only (the CV does not need it:
% the whole document is the credential, and the signature block stays plain).
% Set to {} to drop it. When it is set, the headline should NOT also end in
% \"Ph.D. in Computer Science\" -- one or the other, never both.
\newcommand{\FactPostNominal}{, Ph.D.}
% Use the city you live in. Only write "New York, NY" if that's true.
\newcommand{\FactCity}{NYC metro area}          % confirmed 2026-09-12
% The phone number stays out of this public repo. Define \FactPhone in the
% git-ignored master/contact.local.tex, e.g.
%   \newcommand{\FactPhone}{\slink{tel:+1XXXXXXXXXX}{XXX-XXX-XXXX}}
% Without that file the header simply leaves the phone out.
\IfFileExists{master/contact.local.tex}{\input{master/contact.local.tex}}{}
\newcommand{\FactEmail}{\slink{mailto:carlos@carlosetejada.com}{carlos@carlosetejada.com}}
\newcommand{\FactWebsite}{\slink{https://carlosetejada.com}{carlosetejada.com}}
% Put your LinkedIn handle in BOTH places below, then remove \unverified.
\newcommand{\FactLinkedIn}{\slink{https://www.linkedin.com/in/ctejada10}{linkedin.com/in/ctejada10}}
\newcommand{\FactGitHub}{\slink{https://github.com/ctejada10}{github.com/ctejada10}}
\newcommand{\FactScholar}{\slink{https://scholar.google.com/citations?user=oo_ptugAAAAJ}{Google Scholar}}
% Work status is deliberately left off (decision of 2026-09-11).

% ---------------- University of Washington ----------------
\newcommand{\FactUWOrg}{University of Washington}
% You work remotely from the US; a screener reads the most recent location to
% place you. Inherits \FactCity so there is one city in one place.
\newcommand{\FactUWLocation}{Remote (\FactCity)}
% cv/employment.tex says "Postdoctoral Fellow"; UW's HR title for postdocs is
% usually "Postdoctoral Scholar". Use exactly what your appointment letter says.
\newcommand{\FactUWTitle}{Postdoctoral Researcher}    % confirmed 2026-09-12
% Confirmed 2026-09-11. cv/employment.tex said 2023; corrected there too.
\newcommand{\FactUWDates}{2025 -- Present}
\newcommand{\FactUWStudyN}{12}                  % blind and low-vision participants (confirmed 2026-09-11)
\newcommand{\FactUWPaperVenue}{CHI 2027}        % lead-author paper, submitted Sep 2026, under review. Topic stays OUT of every built PDF: the listing carries authorship position, venue and status only.
% The title stays OFF the resume while the paper is under double-blind review;
% describe it instead. The title itself, and the preflight rule that keeps it
% out of every build, live in master/blocked.txt (untracked).
\newcommand{\FactUWPaperDesc}{Lead author on a multidisciplinary submission}
\newcommand{\FactUWMentees}{four}     % undergraduate researchers you mentor (2026-09-12)
\newcommand{\FactUWTeam}{eight}        % faculty, postdoctoral and undergraduate researchers (2026-09-12)
% Latency improvement on the Bluetooth/HID path, your estimate (2026-09-11).
% Be ready to say what was measured: command -> pins updated, before vs after.
\newcommand{\FactUWLatency}{30\%}

% ---------------- Aarhus University ----------------
\newcommand{\FactAUOrg}{Aarhus University}
\newcommand{\FactAULocation}{Aarhus, Denmark}
% cv/employment.tex says "Postdoctoral Fellow"; Danish contracts usually say "Postdoc".
\newcommand{\FactAUTitle}{Postdoctoral Researcher}    % confirmed 2026-09-12
\newcommand{\FactAUDates}{2023 -- 2024}                % confirmed 2026-09-11
% Neither Aarhus tool was evaluated, so both bullets claim capability and no
% outcome. Any result sentence here needs a study behind it first.

\newcommand{\FactAUProjects}{three}     % projects worked on at Aarhus (confirmed 2026-09-12)

% ---------------- Netcompany ----------------
\newcommand{\FactNCOrg}{Netcompany A/S}
\newcommand{\FactNCShort}{Netcompany}                 % in prose
\newcommand{\FactNCAbout}{listed Danish IT services firm, 7,000+ staff}
\newcommand{\FactNCLocation}{Copenhagen, Denmark}
% Official title per cv/employment.tex, and what an employment check returns.
% Use exactly this; confirm against your contract.
\newcommand{\FactNCTitle}{IT Consultant}
\newcommand{\FactNCDates}{2021 -- 2023}               % Jun 2021 -- Jun 2023; bare years keep the page consistent
\newcommand{\FactNCTeams}{six}
\newcommand{\FactNCLanguages}{four}
\newcommand{\FactNCFone}{87\%}                  % F1. Add the baseline it replaced if you can find it.
\newcommand{\FactNCOverhead}{30\%}
\newcommand{\FactNCToolTeams}{four}
\newcommand{\FactNCClients}{three}

% ---------------- Ph.D. (University of Copenhagen) ----------------
\newcommand{\FactPhDOrg}{University of Copenhagen}
\newcommand{\FactPhDLocation}{Copenhagen, Denmark}
\newcommand{\FactPhDTitle}{Ph.D. Fellow}
\newcommand{\FactPhDDates}{2018 -- 2021}        % employment at UCPH; the Ph.D. began at RIT in 2017
% Figures below come from the papers themselves (PDFs on carlosetejada.com).
% AirTouch and EchoTube report technical evaluations with no participant study
% and no significance tests; Blowhole collected data from ten people. There is
% no counterbalanced, significance-tested study in the record -- see
% master/retired-claims.txt (git-ignored) for what that rules out.
\newcommand{\FactPhDAccuracy}{91--100\%}   % AirTouch, across its four test objects (CHI 2020)
\newcommand{\FactPhDBlowAccuracy}{98\%}    % Blowhole (GI 2018)
\newcommand{\FactPhDSamples}{830}          % Blowhole: "830 blow segments from ten participants"
\newcommand{\FactPhDDataParticipants}{ten}  % Blowhole data collection
\newcommand{\FactPhDTheses}{two}                  % master's theses supervised (cv/supervision.tex)

% ---------------- Rochester Institute of Technology ----------------
\newcommand{\FactRITOrg}{Rochester Institute of Technology}
\newcommand{\FactRITLocation}{Rochester, NY}
% Two appointments, stacked under one employer header (2026-09-11):
% Ph.D. Fellow 2017--2018 (the doctorate's first year, before Copenhagen) and
% Graduate Research Assistant 2014--2016 (alongside the M.Sc.).
\newcommand{\FactRITTitle}{Ph.D. Fellow}
\newcommand{\FactRITDatesPhD}{2017 -- 2018}
\newcommand{\FactRITTitleGRA}{Graduate Research Assistant}
\newcommand{\FactRITDatesGRA}{2014 -- 2016}
% Bitey figures, per the paper (carlosetejada.com/documents/papers/bitey.pdf):
%   78% mean accuracy (SD 15%), above 90% for some participants
%   noise: 1.64 h of recordings, 83,942 candidates, 674 (.8%) false positives,
%          5,255 of 44,324 clicks missed (11.9%); noise is a class in the SVM
\newcommand{\FactBiteyAccuracy}{78\%}
\newcommand{\FactBiteyFP}{0.8\%}
\newcommand{\FactBiteyMissed}{11.9\%}
\newcommand{\FactBiteyN}{20}
% M.Sc. capstone, "Knock on Wood" (May 2016), sole author. Figures from the
% capstone document: SVM 84.03%, random forest 72.92%, naive Bayes 57.64% over
% five materials; the 5-fold cross-validated model scored 97% and then failed
% on new recordings, because samples from one microphone placement fell on
% both sides of the split.
\newcommand{\FactKoWAccuracy}{84\%}
\newcommand{\FactKoWClasses}{five}
\newcommand{\FactKoWLeak}{97\%}
% No count on file for the RIT undergraduates; the bullet works without one.
% If you remember a number, put it here and say "Mentored N undergraduate...".


% ---------------- Tous Software ----------------
\newcommand{\FactTousOrg}{Tous Software Corp.}
\newcommand{\FactTousLocation}{Miramar, FL}
% One line so a reader knows what the company is; the bullets mention telecom
% partners and carrier systems. Confirm the wording.
\newcommand{\FactTousAbout}{CRM and PBX software}     % confirmed 2026-09-12
\newcommand{\FactTousTitle}{Senior Software Developer}
% LinkedIn (written at the time) records Tous as Feb 2012 -- Jul 2014, and a
% Junior Software Engineer role at Synergies Corp. Oct 2011 -- Feb 2012. The
% "2011" belongs to Synergies, not Tous. Industry total is unchanged at 5.
\newcommand{\FactTousDates}{2012 -- 2014}
\newcommand{\FactTousTeam}{three}             % three other developers; the bullet says "led three developers"
\newcommand{\FactTousClients}{five}           % "at least 5", enterprise; stated as five to stay conservative

% ---------------- Synergies (academic CV only; usually omitted) ----------------
\newcommand{\FactSynOrg}{Synergies Strategic Services}
\newcommand{\FactSynLocation}{Santiago, Dominican Republic}
\newcommand{\FactSynTitle}{Junior Software Developer}
\newcommand{\FactSynDates}{2011 -- 2012}

% ---------------- Education ----------------
\newcommand{\FactPhDDegree}{Ph.D., Computer Science}
\newcommand{\FactPhDYear}{2021}
\newcommand{\FactMScDegree}{M.Sc., Information Science}
\newcommand{\FactMScYear}{2016}
\newcommand{\FactBScDegree}{B.Sc., Systems Engineering}
\newcommand{\FactBScOrg}{Pontificia Universidad Católica Madre y Maestra}
\newcommand{\FactBScYear}{2012}    % confirmed 2026-09-12; matches cv/education.tex

% ---------------- Publications ----------------
\newcommand{\FactPubs}{ten}          % peer-reviewed: 8 full papers + CHI'20 doctoral consortium + CHI'20 SIG
\newcommand{\FactPapers}{eight}         % full conference papers
\newcommand{\FactFirstAuthor}{four}    % first-authored full papers: AirLogic, AirTouch, EchoTube, Blowhole
\newcommand{\FactVenues}{CHI, UIST, ISS, TEI, MobileHCI and GI}

% ---------------- Totals used in summaries ----------------
% Recompute once the dates above are settled.
\newcommand{\FactYearsIndustry}{five}   % Synergies 2011-12 + Tous 2012--2014 + Netcompany 2021--2023
\newcommand{\FactYearsResearch}{six}   % Ph.D. 2018--21 (3) + RIT doctoral yr (1) + Aarhus (1) + UW (1+)

\hypersetup{pdftitle={Carlos E. Tejada - Resume}, pdfauthor={Carlos E. Tejada}, pdfsubject={Resume}}

% ---------------- Headline: the line under your name ----------------
% A field label, not a job title you haven't held. Override per application
% with \def\HeadlineOverride{...} if a posting calls the field something else.
\headlinefor{ds}{Data Science \& Applied Machine Learning \textbar{} Experiment Design \& Evaluation}
\headlinefor{mle}{Machine Learning Engineering \textbar{} Python, NLP \& LLM Applications}
\headlinefor{de}{Data Engineering \textbar{} Azure, PySpark \& SQL}
\headlinefor{rs}{Human-Centered AI, Digital Fabrication \& Accessibility}
\headlinefor{uxr}{Mixed-Methods UX Research \textbar{} Accessibility \& AI Tools}

% ---------------- Skills blocks (only what you can defend in an interview) ----------------
% Also on file, but leave off unless you can discuss them: XGBoost, Kubernetes,
% AWS, Databricks, Delta Lake, Azure ML, React, JavaScript, PHP, Matplotlib,
% Seaborn, Autodesk Meshmixer, Linux, CI/CD.
\newcommand{\skillsblock}{%
  \msection{Skills}
  \skillsfor{ds}{%
    \skillline{Methods}{Statistics and experiment design (parameter studies, cross-validation, error analysis) · classification · regression · time series · NLP · retrieval-augmented generation}
    \skillline{Data \& code}{Python (pandas, NumPy, scikit-learn, PyTorch, TensorFlow) · SQL (PostgreSQL, Spark SQL, SQL Server) · PySpark · Azure (Data Lake, Functions)}
    \skillline{Tools}{Power BI · Git · Docker · Flask · FastAPI · OpenAI API}}
  \skillsfor{mle}{%
    \skillline{ML}{PyTorch · TensorFlow · scikit-learn · OpenAI API · retrieval-augmented generation · model monitoring and retraining}
    \skillline{Engineering}{Python · SQL · Java · Go · Flask · FastAPI · REST APIs · Docker · Git}
    \skillline{Data}{PySpark · Spark SQL · PostgreSQL · SQL Server · Azure (Data Lake, Functions)}}
  \skillsfor{de}{%
    \skillline{Pipelines}{PySpark · Spark SQL · Azure (Data Lake, Functions) · ETL · data modeling}
    \skillline{Databases \& code}{PostgreSQL · SQL Server · Python · SQL · Java}
    \skillline{Operations}{Role-based access control · audit logging · Docker · Git · Power BI}}
  \skillsfor{rs}{%
    \skillline{Research methods}{Empirical evaluation of interactive systems · qualitative studies (think-aloud, thematic analysis) · participatory design}
    \skillline{Systems}{Python · machine learning (scikit-learn, PyTorch) · signal processing · LLM APIs · geometry processing (trimesh, Shapely) · OpenSCAD · refreshable tactile displays}
    \skillline{Fabrication}{3D printing · pneumatic, acoustic and ultrasonic sensing}}
  \skillsfor{uxr}{%
    \skillline{Methods}{Think-aloud usability studies · participatory design · thematic analysis · questionnaires · interaction-log analysis}
    \skillline{Analysis}{Python (pandas) · SQL · statistics (descriptive analysis, cross-validation) · data visualization}
    \skillline{Prototyping}{Python/Flask web apps · screen-reader accessibility · tactile displays · 3D printing}}
  \printskills}

\begin{document}
\masterheader

% =====================================================================
\msection{Summary}
% Under 50 words, checkable facts only. Override per application with
% \def\SummaryOverride{...}. Checkable facts only; no career totals.
\summaryfor{ds}{Computer scientist with \FactYearsIndustry{} years in industry and \FactYearsResearch{} in applied-ML research. At \FactNCShort, built Azure/PySpark pipelines over billions of records a month and a \FactNCLanguages-language document classifier (\FactNCFone{} F1); at the University of Washington, built a retrieval-augmented LLM prototype and led a \FactUWStudyN-participant study.}
\summaryfor{mle}{Computer scientist who builds ML systems end to end, from PySpark pipelines on Azure to Python services (Flask, FastAPI) and LLM features. Shipped a \FactNCLanguages-language document classifier (\FactNCFone{} F1) that stayed accurate through data drift, and a real-time system that streams 3D geometry to braille displays.}
\summaryfor{de}{Computer scientist with production data-engineering experience at \FactNCShort: Azure/PySpark ETL over billions of vehicle-telemetry records a month, Spark SQL and PostgreSQL, role-based access control and audit logging across \FactNCClients{} client systems, and self-service tooling that cut reporting overhead \FactNCOverhead.}
\summaryfor{rs}{HCI researcher building interactive systems where machine learning, digital fabrication and accessibility meet. \FactPubs{} peer-reviewed publications at \FactVenues{} (Honorable Mention, ISS 2019). Now leading nonvisual CAD research for blind and low-vision people at the University of Washington, and lead author on a multidisciplinary paper under review at \FactUWPaperVenue.}
\summaryfor{uxr}{Mixed-methods researcher who pairs study design with engineering: builds and instruments the prototype, runs the sessions, and analyzes transcripts and interaction logs together. Most recently led a qualitative study with \FactUWStudyN{} blind and low-vision participants; earlier designed the evaluations behind \FactFirstAuthor{} first-author papers on new sensing techniques.}
\printsummary

% =====================================================================
\msection{Experience}

% ---------------- University of Washington ----------------
% Ownership settled 2026-09-11: the ranking inside YOUR prototype is yours.
% The summer 2026 view-selection / slice-scoring pipelines are your mentees'
% work and are credited as mentoring (uw-mentor), not as your engineering.
\onepagemax{2}
\cvrole{\FactUWOrg}{\FactUWLocation}{\FactUWTitle\ifarch{rs,uxr}{ (HCI \& Accessibility)}{ (Applied ML \& Accessibility)}}{\FactUWDates}
\bul{uw-rag}{ds,mle}{Prototyped a retrieval-augmented LLM companion (Python, OpenAI API): ranks a model's cross-sections, feeds the top-k most informative to the LLM, and returns descriptions of shapes and of changes between versions.}
\bul{uw-rag.rs}{rs}{Built a retrieval-augmented LLM companion for 3D captioning: ranks every cross-section of a model, retrieves the top-k most informative slices, and grounds an LLM (OpenAI API) in them to describe shapes and changes between versions.}
\bul{uw-study.ds}{ds}{Designed and ran an IRB-approved study with \FactUWStudyN{} blind and low-vision participants, analyzing interaction logs in Python alongside session transcripts; lead author on the resulting submission.}
\bul{uw-study}{rs,uxr}{Led an IRB-approved qualitative study with \FactUWStudyN{} blind and low-vision participants, who compared 3D-printed and virtual models while thinking aloud; lead author on the resulting submission.}
\bul{uw-system.ds}{ds,mle,de}{Lead developer of a nonvisual CAD system: a Python backend slices 3D models on demand and streams each cross-section to refreshable braille displays from a browser app that needs no install; released as open source.}
\bul{uw-system}{rs,uxr}{Lead developer of a nonvisual CAD system that lets blind and low-vision people explore and edit 3D models by touch: a Python backend slices STL geometry on demand and streams it to refreshable pin displays (APH Monarch, Dot Pad) from a screen-reader-accessible web app.}
\bul{uw-logs}{de,uxr}{Instrumented every command in the system and analyzed participants' interaction logs in Python (per-task timelines, command frequencies) alongside thematic analysis of the session transcripts.}
\bul{uw-codesign}{rs,uxr}{Designed the interaction model with a DeafBlind co-designer and co-author through sustained participatory design: axis-aligned slicing, fixed-step depth control, and outline, fill and cut rendering modes.}
\bul{uw-infra}{mle,de}{Cut latency to the displays by \FactUWLatency{} by rebuilding the Bluetooth/HID path, and moved the backend onto lab servers so participants run everything in Chrome with nothing installed.}
\bul{uw-wizard}{uxr}{Built a guided, in-app study flow that walks participants through models and questionnaires, keeping every session consistent.}
\bul{uw-mentor}{mle,rs}{Set technical direction for an \FactUWTeam-person team of faculty, postdoctoral and undergraduate researchers, and mentor \FactUWMentees{} undergraduates building view-selection and slice-scoring pipelines.}
\bul{uw-workshop}{uxr}{Co-facilitated a tactile design workshop with blind and low-vision students, and demoed the system at an assistive-technology meetup at the New York Public Library.}
% The LLM descriptions have had no formal evaluation (confirmed 2026-09-11), so
% uw-rag claims no quality result -- "prototyped" is the honest verb. If one ever
% runs (N descriptions rated, error categories, grounded vs ungrounded), add it
% to uw-rag: that sentence is the single biggest upgrade available to this role.
\printbullets

% ---------------- Aarhus University ----------------
% Settled 2026-09-11: neither tool was evaluated and neither was demoed
% publicly, so both bullets describe capability and claim no outcome. The
% time-series work was early-stage, so the verbs say so. The third project
% produced a paper you are not an author on -- not resume material.
\onepagemax{1}
\cvrole{\FactAUOrg}{\FactAULocation}{\FactAUTitle\ifarch{rs,uxr}{ (LLMs \& Digital Fabrication)}{ (Applied ML \& LLM Systems)}}{\FactAUDates}
\bul{aarhus-llm-cad}{ds,mle,rs}{Built an LLM-powered design tool (Python, FastAPI, OpenSCAD) that turns a plain-language prompt into printable CAD for interactive devices, validating the generated geometry against fabrication constraints so what reaches the printer is buildable.}
% If the MODEL emitted the OpenSCAD (rather than you writing it), swap in this
% wording -- it is the difference between a toolchain bullet and a code-generation
% one, and it answers CodeLLM experience on AI-for-code postings:
%\bul{aarhus-llm-cad.code}{mle,rs}{Built an LLM-powered design tool (Python, FastAPI) that writes OpenSCAD from a plain-language prompt and checks the generated code's geometry against fabrication constraints before anything reaches the printer.}
\bul{aarhus-nlq}{ds,mle,de,uxr}{Built a natural-language query layer (Python, FastAPI, LLM) that let non-technical collaborators pull answers from research datasets without writing queries.}
\bul{aarhus-ts}{ds,mle,de,rs}{Prototyped and benchmarked time-series models (PyTorch, scikit-learn) that turn sensor signals from 3D-printed devices into discrete events, comparing candidates against baselines on accuracy, latency and reliability to set the modeling direction for \FactAUProjects{} projects.}
\printbullets

% ---------------- Netcompany ----------------
% The domain is vehicle telemetry for traffic analytics (free-flow speed,
% congestion). Wording below reconstructs that from your description on
% 2026-09-11 -- correct any term that isn't how the work was actually framed.
\onepagemax{2}
\cvrole{\FactNCOrg\ {\mdseries\color{text}(\FactNCAbout)}}{\FactNCLocation}{\FactNCTitle\ (Data \& ML)}{\FactNCDates}
\bul{nc-nlp}{ds,mle,rs}{Replaced manual triage at a government client with a \FactNCLanguages-language document classifier (scikit-learn, TensorFlow) that routes tens of thousands of documents a month to the right team: \FactNCFone{} F1, held through drift with periodic retraining.}
\bul{nc-etl}{ds,mle,de,rs,uxr}{Built an Azure/PySpark ETL pipeline (Spark SQL, PostgreSQL) that turns billions of vehicle-telemetry records a month into trip-level tables behind traffic speed and travel-time models; technical contact for \FactNCTeams{} sales, operations and business-development teams.}
\bul{nc-tools}{ds,de,uxr}{Cut reporting overhead \FactNCOverhead{} for \FactNCToolTeams{} teams by automating data workflows (Python, SQL, REST APIs, Azure Functions) and building Power BI dashboards on free-flow speed and congestion that senior leadership used weekly.}
\bul{nc-sparktuning}{de,mle}{Profiled Spark SQL execution plans and tuned partitioning and broadcast joins to cut job runtimes.}
\bul{nc-automl}{mle}{Built the retraining loop behind the classifier: benchmarks and tunes candidate models against incoming data as distributions shift, replacing manual experimentation.}
\bul{nc-gov}{de}{Implemented role-based access control and audit logging for sensitive data across \FactNCClients{} client systems to meet contractual and regulatory requirements.}
\printbullets

% ---------------- Ph.D. ----------------
\onepagemax{1}
\cvrole{\FactPhDOrg}{\FactPhDLocation}{\FactPhDTitle\ (Doctoral Researcher, HCI)}{\FactPhDDates}
\bul{phd-experiments}{ds,rs,uxr}{Designed and ran the empirical evaluations behind the \FactFirstAuthor{} first-author papers of my doctorate: parameter sweeps over chamber geometry, outlet spacing and tube dimensions, measured against ground truth, that set the final design of each sensing technique.}
\bul{phd-ml}{ds,mle,de}{Built real-time classification pipelines (Python, scikit-learn; windowed features, RBF-kernel SVMs) for touch on 3D-printed objects (\FactPhDAccuracy{} accuracy, no per-user calibration) and for blow-activated input (\FactPhDBlowAccuracy{}, cross-validated on \FactPhDSamples{} labeled samples from \FactPhDDataParticipants{} participants).}
\bul{phd-tooling}{mle,rs}{Built design tooling on Autodesk Meshmixer's scripting API that generates sensor geometry inside users' own 3D models, so non-experts could fabricate custom interactive objects.}
\bul{phd-mentor}{rs,uxr}{Supervised \FactPhDTheses{} master's theses; one became a TEI 2022 paper (ClipWidgets).}
\printbullets

% ---------------- RIT ----------------
\unlessonepage{%
\onepagemax{1}
\cvrole{\FactRITOrg}{\FactRITLocation}{\FactRITTitle}{\FactRITDatesPhD}
\unlessonepage{\cvsubrole{\FactRITTitleGRA}{\FactRITDatesGRA}}
\bul{rit-acoustic}{ds,mle,de}{Built the acoustic classification pipeline (Python, scikit-learn; FFT features, RBF-kernel SVM) behind Bitey, a tooth-click input technique sensed through a nine-dollar bone-conduction microphone, and ran the \FactBiteyN-participant data collection and evaluation behind it: \FactBiteyAccuracy{} mean accuracy across up to five click classes.}
\bul{rit-noise}{none}{Made it usable outside the lab by training a noise class from 1.6 hours of talking, walking, eating and driving: \FactBiteyFP{} of 84,000 non-click events misclassified as clicks, at the cost of \FactBiteyMissed{} missed clicks.}
\bul{rit-capstone}{ds,mle}{Built an acoustic material classifier for my M.Sc.\ capstone (Python, scikit-learn; FFT features over a swept tone) so laser cutters and CNC mills could load the right settings automatically: \FactKoWAccuracy{} across \FactKoWClasses{} materials, more than ten points above random forest and naive Bayes.}
\bul{rit-leakage}{ds,mle,de}{Found the leak behind it: five-fold cross-validation scored \FactKoWLeak{} because recordings from a single microphone placement landed on both sides of the split; rebuilt the evaluation on held-out placements, which is where the \FactKoWAccuracy{} comes from.}
\bul{rit-apps}{uxr}{Built two live applications on the tooth-click technique, a list selector and an on-screen keyboard, to test it in use rather than on recordings.}
\bul{rit-mentor}{ds,mle,rs,uxr}{Mentored undergraduate researchers through their lab projects, from study design and data collection to feature extraction, classifier training and analysis.}
\bul{rit-bitey}{rs,uxr}{Co-authored Bitey (MobileHCI 2016), hands-free interface control through tooth-click gestures, with clinical researchers at the University of Rochester Medical Center.}
\printbullets}

% ---------------- Tous Software ----------------
% Printed for industry builds only (2026-09-11). On a research or UX resume a
% 2012--14 developer job is filler, and dropping it costs no timeline: RIT
% picks up in 2014. Any application can re-enable it from the bank.
\unlessonepage{\onlyfor{ds,mle,de}{%
\cvrole{\FactTousOrg\ {\mdseries\color{text}(\FactTousAbout)}}{\FactTousLocation}{\FactTousTitle}{\FactTousDates}
\bul{tous-portal}{ds,mle,de}{Led \FactTousTeam{} developers building the customer-facing reporting portal across the company's CRM and PBX products (Java, PHP, Go), plus the REST APIs telecom partners integrated with; \FactTousClients{} enterprise clients used it, their leadership running the reports for operating decisions.}
\bul{tous-integrations}{de,mle}{Replaced manual coordination between carriers with automated data exchange, building the inbound integrations to their systems.}
\bul{tous-db}{de}{Designed the database schema and tuned the queries behind clients' on-demand reports across both product lines.}
\printbullets}}

% ---------------- Synergies (catalog only) ----------------
\catalogonly{%
  \cvrole{\FactSynOrg}{\FactSynLocation}{\FactSynTitle}{\FactSynDates}
  {\fontsize{9pt}{11pt}\bodyfont\color{gray}From the academic CV; no bullets on file. Usually omitted.}\par}

% =====================================================================
% Skills (every role type except rs puts them here; rs puts them after publications)
\ifcatalog\skillsblock\else\ifarch{rs}{}{\skillsblock}\fi

% =====================================================================
% Publications
\newcommand{\pubsummary}{\FactPubs{} peer-reviewed publications, including \FactPapers{} full papers at \FactVenues{} (\FactFirstAuthor{} first-authored); Honorable Mention, ISS 2019. Full list: \FactScholar.}
\newcommand{\pubunderreview}{\textit{Under review:} \FactUWPaperDesc{} to \FactUWPaperVenue. Title and details withheld during double-blind review.}
\newcommand{\pubAirTouch}{\me, R. Ramakers, S. Boring, D. Ashbrook. \href{https://doi.org/10.1145/3313831.3376136}{AirTouch: 3D-printed Touch-Sensitive Objects Using Pneumatic Sensing}. CHI 2020.}
\newcommand{\pubEchoTube}{\me, J. McIntosh, K. A. Bergen, S. Boring, D. Ashbrook, A. Marzo. \href{https://doi.org/10.1145/3343055.3359712}{EchoTube: Robust Touch Sensing along Flexible Tubes using Waveguided Ultrasound}. ISS 2019. \textbf{Honorable Mention.}}
\newcommand{\pubAirLogic}{\me, V. Savage, M. Zhong, R. Ramakers, D. Ashbrook, H. Kim. \href{https://doi.org/10.1145/3526113.3545642}{AirLogic: Embedding Pneumatic Computation and I/O in 3D Models to Fabricate Electronics-Free Interactive Objects}. UIST 2022.}
\newcommand{\pubLaCir}{N. C. Buch, \me, V. Savage, D. Ashbrook. \href{https://doi.org/10.1145/3613904.3642888}{LaCir: A Multilayered Laser-cuttable Material to Co-fabricate Circuitry and Structural Components}. CHI 2024.}
\newcommand{\pubClipWidgets}{A. Visschedijk, H. Kim, \me, D. Ashbrook. \href{https://doi.org/10.1145/3490149.3501314}{ClipWidgets: 3D-printed Modular Tangible UI Extensions for Smartphones}. TEI 2022.}
\newcommand{\pubMorpheesPlug}{H. Kim, A. Everitt, \me, M. Zhong, D. Ashbrook. MorpheesPlug: A Toolkit for Prototyping Shape-Changing Interfaces. CHI 2021.}
\newcommand{\pubBlowhole}{\me, O. Fujimoto, Z. Li, D. Ashbrook. \href{https://doi.org/10.20380/GI2018.18}{Blowhole: Blowing-Activated Tags for Interactive 3D-Printed Models}. GI 2018.}
\newcommand{\pubBitey}{D. Ashbrook, \me, D. Mehta, A. Jiminez, G. Muralitharam, S. Gajendra, R. Tallents. \href{https://doi.org/10.1145/2935334.2935389}{Bitey: An Exploration of Tooth Click Gestures for Hands-Free User Interface Control}. MobileHCI 2016.}
\newcommand{\pubOther}{\textit{Also peer-reviewed:} \me. Print-and-Play: 3D-printed Interactive Objects Without Assembly or Calibration. CHI 2020 Extended Abstracts (Doctoral Consortium). \textbullet{} A. Alvarado Garcia, \dots, \me, M. Wong-Villacres. Fostering HCI Research in, by, and for Latin America. CHI 2020 Extended Abstracts (SIG).}

\ifcatalog
  \msection{Publications}
  {\fontsize{8.5pt}{11pt}\bodyfont\color{gray}Catalog: rs prints the full list below. ds and mle print the summary line + AirTouch, EchoTube, AirLogic. uxr prints the summary line + under-review paper, AirTouch, Bitey. de prints the summary line only.}\par\vspace{3pt}
  \pubitem{\pubsummary}\pubitem{\pubunderreview}\pubitem{\pubLaCir}\pubitem{\pubAirLogic}\pubitem{\pubClipWidgets}\pubitem{\pubMorpheesPlug}\pubitem{\pubAirTouch}\pubitem{\pubEchoTube}\pubitem{\pubBlowhole}\pubitem{\pubBitey}\pubitem{\pubOther}
\else\ifonepage
  \msection{Publications}\pubitem{\pubsummary}
\else
  \ifarch{ds,mle}{\msection{Selected Publications}\pubitem{\pubsummary}\pubitem{\pubAirTouch}\pubitem{\pubEchoTube}\pubitem{\pubAirLogic}}{}
  \ifarch{uxr}{\msection{Selected Publications}\pubitem{\pubsummary}\pubitem{\pubunderreview}\pubitem{\pubAirTouch}\pubitem{\pubBitey}}{}
  \ifarch{de}{\msection{Publications}\pubitem{\pubsummary}}{}
  \ifarch{rs}{\msection{Publications}\pubitem{\pubunderreview}\pubitem{\pubLaCir}\pubitem{\pubAirLogic}\pubitem{\pubClipWidgets}\pubitem{\pubMorpheesPlug}\pubitem{\pubAirTouch}\pubitem{\pubEchoTube}\pubitem{\pubBlowhole}\pubitem{\pubBitey}\pubitem{\pubOther}}{}
\fi\fi

% Skills for research-scientist builds
\ifcatalog\else\ifarch{rs}{\skillsblock}{}\fi

% =====================================================================
% Service, talks and press (research scientist builds and the catalog)
\unlessonepage{\onlyfor{rs}{%
  \msection{Service \& Talks}
  \pubitem{Program committee, Late-Breaking Work, CHI 2021 \textbullet{} Assistant to the Program Committee Chairs, Accessibility and Aging, CHI 2020.}
  \pubitem{Reviewer: CHI, UIST, TOCHI, IMWUT, ISS, TEI, MobileHCI, IEEE Pervasive Computing (2018--2023).}
  \pubitem{Invited talk: ``Enabling Non-experts to Author Tangible Interactions,'' Human-Computer Integration Lab, Chicago, 2020 \textbullet{} Press: Arduino Blog on AirTouch, 2020.}}}
\catalogonly{%
  \pubitem{\textit{Training (catalog only):} UIST School, Hasso Plattner Institute, 2018 \textbullet{} ACM SIGCHI Summer School on Computational Interaction, ETH Zurich, 2017.}}

% =====================================================================
\msection{Education}
\eduline{\FactPhDDegree}{\FactPhDOrg}{\FactPhDYear}
\eduline{\FactMScDegree}{\FactRITOrg}{\FactMScYear}
\eduline{\FactBScDegree}{\FactBScOrg}{\FactBScYear}

\end{document}

% Retired claims (what was removed and why) live in master/retired-claims.txt,
% which is git-ignored: this repository is public.
